% Rev 9 -- Targeted Embedded Controls / Turf & Compact Utility Software
\documentclass[10pt,letterpaper]{article}

\usepackage[margin=0.50in]{geometry}
\usepackage[T1]{fontenc}
\usepackage[utf8]{inputenc}
\usepackage[default,type1]{sourcesanspro}
\usepackage[final]{microtype}
\usepackage[explicit]{titlesec}
\usepackage{enumitem}
\usepackage{tabularx}
\usepackage{array}
\usepackage{ragged2e}
\usepackage{iftex}
\usepackage[hidelinks]{hyperref}

\ifPDFTeX
  \input{glyphtounicode}
  \pdfgentounicode=1
\fi

\hypersetup{
  pdftitle={Connor Savugot Resume},
  pdfauthor={Connor Savugot},
  pdfsubject={Embedded Controls Software Engineering -- MATLAB, C, CAN J1939},
  pdfkeywords={embedded software, embedded C, C++, MATLAB, dynamic systems, transfer functions, PID control, controller tuning, CAN, J1939, FreeRTOS, real time systems, functional testing, software requirements, component integration, CARLA, Lustre, robotics, autonomous vehicle, dsPIC33CK, PIC32, TMS320 DSP, STM32G474, JTAG, oscilloscopes, logic analyzers, Git}
}

\pagestyle{empty}
\setcounter{secnumdepth}{0}
\setlength{\parindent}{0pt}
\setlength{\parskip}{0pt}
\hfuzz=0.4pt

\titleformat{\section}
  {\fontsize{11.3}{11.9}\selectfont\bfseries\uppercase}
  {}{0pt}{#1\hspace{0.12cm}\titlerule[0.65pt]}
\titlespacing{\section}{0pt}{0.14cm}{0.05cm}

\newlist{resumebullets}{itemize}{1}
\setlist[resumebullets]{
  label=\textbullet,
  leftmargin=0.42cm,
  topsep=0.004cm,
  itemsep=0.004cm,
  parsep=0pt,
  partopsep=0pt
}

\newcommand{\companyheader}[1]{%
  {\fontsize{10.4}{10.8}\selectfont\bfseries #1}\par\vspace{-0.03cm}
}

\newcommand{\resumeentry}[2]{%
  \begin{tabularx}{\textwidth}{@{}X>{\raggedleft\arraybackslash}p{2.40in}@{}}
    \textbf{#1} & #2
  \end{tabularx}\vspace{-0.075cm}
}

\newcommand{\roleentry}[2]{%
  \begin{tabularx}{\textwidth}{@{}X>{\raggedleft\arraybackslash}p{2.40in}@{}}
    \textbf{#1} & #2
  \end{tabularx}\vspace{-0.075cm}
}

\newcommand{\descriptorentry}[1]{%
  \begin{tabularx}{\textwidth}{@{}X>{\raggedleft\arraybackslash}p{2.40in}@{}}
    \textit{#1} &
  \end{tabularx}\vspace{-0.075cm}
}

\begin{document}
\fontsize{10.00}{11.00}\selectfont
\setlength{\RaggedRightRightskip}{0pt plus 5em}
\RaggedRight
\emergencystretch=1em
\hyphenpenalty=10000
\exhyphenpenalty=10000

\begin{center}
  {\fontsize{21.5}{22.2}\selectfont\bfseries Connor Savugot}\\[1.5pt]
  {\bfseries B.S. in Computer Engineering | North Carolina State University | Graduated May 2026}\\[1.5pt]
  {\fontsize{9.6}{10.2}\selectfont
    \href{mailto:consav04@gmail.com}{consav04@gmail.com}
    \enspace|\enspace \href{tel:+18285419487}{+1-828-541-9487}
    \enspace|\enspace Murphy, NC
    \enspace|\enspace \href{https://linkedin.com/in/connorsavugot}{linkedin.com/in/connorsavugot}
    \enspace|\enspace \href{https://github.com/Clem085}{github.com/Clem085}}
\end{center}
\vspace{-0.18cm}

\section{Embedded Controls Profile}
Embedded firmware and controls engineer developing C software for power and communications systems. Combines professional MATLAB data analysis, CAN/J1939, FreeRTOS, and hardware integration with academic PID, dynamic-system modeling, robotics, and autonomous-vehicle controls.

\section{Relevant Skills}
\textbf{Languages \& Engineering Analysis:} Embedded C, C++, MATLAB, Python, Lustre\\
\textbf{Controls \& Real-Time:} Dynamic-system modeling, PID/PI/PD tuning, cascaded and discrete-time control, state-space and transfer-function analysis\\
\textbf{Platforms \& Operating Systems:} FreeRTOS, bare metal, Embedded Linux; dsPIC33CK, PIC32, TMS320 DSP, STM32G474; interrupts/timers, peripheral/HAL driver bring-up\\
\textbf{Vehicle Networks \& I/O:} Classical CAN (CAN 2.0), CAN FD, J1939; ADC, PWM, I2C, SPI, UART/RS-232\\
\textbf{Development, Delivery \& Validation:} Requirements analysis, interface specifications, software integration, unit, component, and functional testing, design reviews, Git/GitLab, GCC/Make build automation, Code Composer Studio, JTAG, oscilloscopes, logic/protocol analyzers, schematics/datasheets

\section{Engineering Experience}
\companyheader{AEGIS Power Systems}
\roleentry{Embedded Firmware Engineer}{May 2026--Present}
\begin{resumebullets}
  \item Built a MATLAB workflow to parse captured CAN telemetry and compare candidate digital filters by reading variation under steady input and time to follow large electrical changes; selected rolling-average and EMA approaches, derived per-channel scaling, and implemented rolling-sum averaging in embedded C.
  \item Develop low-level C firmware for 16-bit dsPIC33CK DSCs and bring up peripherals from schematics, datasheets, reference manuals, and protocol specifications; isolate register, timing, configuration, and electrical faults with JTAG, targeted tests, oscilloscopes, and logic analyzers.
  \item Integrate CAN-connected power hardware with TI AM62 Embedded Linux software, tracing defects across wiring, bus traffic, device configuration, services, and application code; coordinate priorities, findings, and validation through standups with a distributed firmware team.
\end{resumebullets}

\roleentry{Embedded Systems Intern}{May--Aug 2025}
\roleentry{Computer Science Intern}{May--Aug 2024}
\descriptorentry{Selected work across both internships}
\begin{resumebullets}
  \item Developed bare-metal C power-control firmware for a TMS320 DSP driving transformer-coupled bidirectional buck--boost supplies; coordinated switching across 12 PWM A/B pairs (24 outputs) and validated ADC/PWM behavior in Code Composer Studio.
  \item Developed embedded C firmware for a PIC32 running FreeRTOS, coordinating CAN/J1939, I2C, SPI, UART, and update behavior within a two-processor PIC32/TMS320 product; built a multistage update workflow using C++, Python, and embedded firmware to transfer, program, verify, and recover firmware images.
\end{resumebullets}

\companyheader{TEAM Industries}
\roleentry{Engineering Intern}{Jun--Aug 2022 and Jun--Aug 2021}
\begin{resumebullets}
  \item Built a repeatable Excel/VBA dimensional and spline-tolerance validation workflow that flagged out-of-spec parts; maintained structured quality records and collaborated with engineering and manufacturing teams on test procedures and revision tracking.
\end{resumebullets}

\section{Controls \& Vehicle Software Projects}
\resumeentry{Dynamic Systems \& Signal Analysis | MATLAB}{Spring 2024}
\begin{resumebullets}
  \item Built vectorized MATLAB scripts to evaluate an RLC transfer function across frequency and model first-order circuit and second-order mechanical responses, using plots to analyze damping, transient response, and steady-state behavior.
\end{resumebullets}

\resumeentry{Robotics \& Digital Controls | Python}{Fall 2025}
\begin{resumebullets}
  \item Extended Python simulations for mass-spring-damper, planar-VTOL, and Hummingbird models; derived linearized/state-space models and tuned cascaded PD/PID controllers with saturation, anti-windup, and PWM actuator mixing.
\end{resumebullets}

\resumeentry{Embedded Systems for Autonomous Driving | Lustre, Python/CARLA}{Spring 2026}
\begin{resumebullets}
  \item Implemented and tuned a deterministic PID cruise controller in Lustre with initialized state, discrete-time integral/derivative terms, and output limits; compared its simulated response with an earlier controller. Separately used Luke counterexample traces to correct temporal monitors and verify state-machine properties.
  \item On a three-person team, integrated a Python agent with CARLA's planning stack for PID-based longitudinal control; tuned target-speed response and evaluated Town10HD runs using completion time, collisions, and lane invasions.
\end{resumebullets}

\resumeentry{EV Active Sensor Adapter | Embedded C, CAN, Requirements \& Test}{Senior Design}
\begin{resumebullets}
  \item Translated battery-sensing and CAN constraints into system/interface requirements and validation criteria for a team-built STM32G474/BQ79616 adapter; supported design reviews and staged integration testing, including seven passing host-side tests for conversions, CAN payloads, boundary cases, and checksums.
\end{resumebullets}

\end{document}
